% =========================================================================
% Supplementary Annexure: Cross-Domain Benchmark Validation of KI-SSDM
% Companion document to:
% "Knowledge-informed state-space degradation modelling with conformal
%  uncertainty quantification for ball-screw remaining useful life prediction"
% =========================================================================
\documentclass[11pt,a4paper]{article}

\usepackage[margin=2.0cm]{geometry}
\usepackage{amsmath,amssymb,amsfonts,mathtools,bm}
\usepackage{booktabs,tabularx,multirow,array}
\usepackage{graphicx}
\usepackage{xcolor}
\usepackage{subcaption}
\usepackage{enumitem}
\usepackage[colorlinks=true,linkcolor=blue,citecolor=blue,urlcolor=blue]{hyperref}

\newcommand{\R}{\mathbb{R}}
\newcommand{\RUL}{\mathrm{RUL}}
\newcommand{\bPhi}{\bm{\Phi}}
\newcommand{\btheta}{\bm{\theta}}
\newcommand{\Dcrit}{D_{\mathrm{crit}}}

\title{\textbf{Supplementary Annexure: Cross-Domain Benchmark Validation of Knowledge-Informed State-Space Degradation Modelling (KI-SSDM) on the IEEE PHM 2012 / PRONOSTIA Platform}}
\author{\textbf{Companion Document to:}\\
``Knowledge-informed state-space degradation modelling with conformal uncertainty quantification\\
for ball-screw remaining useful life prediction''}
\date{}

\begin{document}
\maketitle

\begin{abstract}
This supplementary annexure provides an independent, authentic, and reproducible cross-domain benchmark validation of the \textbf{Knowledge-Informed State-Space Degradation Model (KI-SSDM)} on the public IEEE PHM 2012 / PRONOSTIA bearing degradation dataset (FEMTO-ST Institute). While the primary manuscript focuses on mechatronic ball-screw feed drives operating under controller-internal telemetry without external accelerometers (avoiding costly run-to-failure testing through physics-informed digital twins and accelerated life testing), this annexure evaluates the mathematical generalizability of the state-space formulation on run-to-failure vibration benchmarks. KI-SSDM is evaluated under two strict zero-leakage protocols: Protocol~A (the official competition challenge split of 6 training and 11 held-out test bearings) and Protocol~B (full leave-one-bearing-out cross-validation across all 17 bearings evaluated at 8 inspection origins, $N=136$ evaluations). Published comparators include: (1)~Bayesian exponential degradation (Gebraeel et al., 2005); (2)~Particle Filter exponential degradation with $3\sigma$ First Predicting Time (Li et al., 2015); (3)~Deep LSTM life-percentage sequence network (Guo et al., 2017; Zhu et al., 2019); alongside two signal-free reference baselines (a survival-conditioned life prior and a constant median-RUL predictor). Under leave-one-bearing-out evaluation, KI-SSDM achieves an MAE of $3{,}828$~s ($3{,}099$~s post-onset), significantly outperforming the Bayesian exponential baseline ($6{,}696$~s, gain $+2{,}868$~s, Holm $p_t = 6 \times 10^{-5}$) and Particle Filter baseline ($7{,}201$~s, gain $+3{,}373$~s, Holm $p_t = 0.0013$). However, KI-SSDM achieves statistical parity with the deep LSTM ($4{,}811$~s, gain $+983$~s, Holm $p_t = 0.50$) and ties the signal-free constant median-RUL predictor ($3{,}744$~s, difference $-84$~s, Holm $p_t = 0.67$). We transparently disclose that on PRONOSTIA, late and abrupt defect onset ($3\sigma$ FPT occurring after 95\% of operating life for 9 of the 17 bearings) combined with severe failure-threshold dispersion ($1.1\times\text{--}31.7\times$ baseline RMS) causes all vibration-based prognostic methods to extract limited information beyond the population survival distribution. Split-conformal 90\% prediction intervals achieve coverage close to nominal ($88.2\%$, within finite-sample binomial bounds), but reflect this underlying physical dispersion through wide uncertainty bands ($\approx 31{,}400$~s).
\end{abstract}

\vspace{0.6em}
\noindent\textbf{Keywords:} Prognostics and Health Management, PRONOSTIA Bearing Dataset, Knowledge-Informed State-Space Model, Split-Conformal Prediction, Accelerated Life Testing.

\section{Introduction and Industrial Context}
\label{sec:annex_intro}

In precision machine tool informatics, remaining useful life (RUL) prognosis is subject to two operational realities:
\begin{enumerate}[label=(\alph*),leftmargin=2.2em]
    \item \textbf{Sensorless Telemetry Ingestion}: High-volume CNC machine tools operate in hostile machining environments with cutting fluid immersion, chips, and inaccessible rotating shafts. External vibration accelerometers require specialized cabling and charge amplifiers that are rarely tolerated in production. Prognostics must instead ingest controller-resident telemetry (drive torque, current, speed, position) via standard OPC~UA interfaces.
    \item \textbf{Avoidance of Run-to-Failure Testing in Production}: In high-value machine tools, running structural drive axes to catastrophic seizure or ball spalling is economically unacceptable. While dedicated laboratory test benches use accelerated life testing (ALT) to failure under severe preloads specifically to calibrate the underlying physical digital twins and empirical wear kinetics, operational production machines avoid destructive failure by monitoring degradation continuously from controller-internal telemetry.
\end{enumerate}

To establish the algorithmic behavior of the proposed state-space formulation on standard benchmark datasets, this annexure evaluates the framework on the public \textbf{PRONOSTIA bearing platform} (PHM benchmark data~\cite{Lei_2018}). Rolling-element bearings share the same Hertzian contact mechanics, Lundberg-Palmgren fatigue laws, and stochastic spalling progressions as ball screws, providing an ideal cross-domain testbed.

\section{PRONOSTIA Experimental Platform and Dataset Characterisation}
\label{sec:annex_setup}

The PRONOSTIA platform comprises 17 miniature ball bearings tested to failure across three operational regimes:
\begin{itemize}[leftmargin=2.2em]
    \item \textbf{Condition 1}: $1{,}800$~rpm ($30.0$~Hz rotational speed), $4{,}000$~N radial load (7 bearings).
    \item \textbf{Condition 2}: $1{,}650$~rpm ($27.5$~Hz rotational speed), $4{,}200$~N radial load (7 bearings).
    \item \textbf{Condition 3}: $1{,}500$~rpm ($25.0$~Hz rotational speed), $5{,}000$~N radial load (3 bearings).
\end{itemize}
Vibration was sampled at $25.6$~kHz for $0.1$~s every $10$~s ($\Delta t = 10$~s per snapshot).

Table~\ref{tab:pronostia_dataset} summarizes the lifetime characteristics, onset timing, and failure thresholds across all 17 bearings.

\begin{table*}[!t]
\centering
\caption{Summary characteristics of all 17 run-to-failure bearings in the IEEE PHM 2012 / PRONOSTIA dataset. Truncation and ground-truth values correspond to the official challenge test protocol.}
\label{tab:pronostia_dataset}
\small
\begin{tabular*}{\textwidth}{@{\extracolsep{\fill}}l c c c c c c c@{}}
\toprule
\textbf{Bearing} & \textbf{Operating} & \textbf{Total Life} & \textbf{Onset FPT} & \textbf{Failure RMS /} & \textbf{Failure RMS} & \textbf{Challenge} & \textbf{Official RUL} \\
\textbf{Identifier} & \textbf{Condition} & \textbf{(s)} & \textbf{(\% of Life)} & \textbf{Baseline} & \textbf{(g)} & \textbf{Truncation (s)} & \textbf{(s)} \\
\midrule
\texttt{Bearing1\_1} & 1 & $28{,}030$ & 66.8\% & $12.1\times$ & 6.71 & --- & --- \\
\texttt{Bearing1\_2} & 1 & $8{,}710$  & 95.2\% & $5.8\times$  & 2.23 & --- & --- \\
\texttt{Bearing2\_1} & 2 & $9{,}110$  & 96.6\% & $5.3\times$  & 2.70 & --- & --- \\
\texttt{Bearing2\_2} & 2 & $7{,}970$  & 30.2\% & $5.9\times$  & 1.97 & --- & --- \\
\texttt{Bearing3\_1} & 3 & $5{,}150$  & 95.9\% & $1.9\times$  & 0.87 & --- & --- \\
\texttt{Bearing3\_2} & 3 & $16{,}370$ & 98.2\% & $5.8\times$  & 2.37 & --- & --- \\
\midrule
\texttt{Bearing1\_3} & 1 & $23{,}750$ & 57.6\% & $26.4\times$ & 10.34 & $18{,}020$ & $5{,}730$ \\
\texttt{Bearing1\_4} & 1 & $14{,}280$ & 4.1\%  & $26.9\times$ & 11.10 & $11{,}390$ & 339 ($2{,}890$)$^\dagger$ \\
\texttt{Bearing1\_5} & 1 & $24{,}630$ & 99.1\% & $3.9\times$  & 1.40  & $23{,}020$ & $1{,}610$ \\
\texttt{Bearing1\_6} & 1 & $24{,}480$ & 66.7\% & $2.9\times$  & 1.29  & $23{,}020$ & $1{,}460$ \\
\texttt{Bearing1\_7} & 1 & $22{,}590$ & 98.0\% & $5.2\times$  & 2.40  & $15{,}020$ & $7{,}570$ \\
\texttt{Bearing2\_3} & 2 & $19{,}550$ & 13.7\% & $31.7\times$ & 14.25 & $12{,}020$ & $7{,}530$ \\
\texttt{Bearing2\_4} & 2 & $7{,}510$  & 99.1\% & $4.4\times$  & 1.56  & $6{,}120$  & $1{,}390$ \\
\texttt{Bearing2\_5} & 2 & $23{,}110$ & Latent & $1.1\times$  & 0.73  & $20{,}020$ & $3{,}090$ \\
\texttt{Bearing2\_6} & 2 & $7{,}010$  & 98.1\% & $4.9\times$  & 1.50  & $5{,}720$  & $1{,}290$ \\
\texttt{Bearing2\_7} & 2 & $2{,}300$  & 97.4\% & $16.8\times$ & 6.28  & $1{,}720$  & 580 \\
\texttt{Bearing3\_3} & 3 & $4{,}340$  & 76.7\% & $4.3\times$  & 1.30  & $3{,}520$  & 820 \\
\bottomrule
\end{tabular*}
\vspace{0.2em}
\raggedright\footnotesize $^\dagger$ Note: \texttt{Bearing1\_4} official challenge key is $339$~s, whereas the full run-to-failure file record gives $2{,}890$~s. Both evaluations are reported.
\end{table*}

Inspection of Table~\ref{tab:pronostia_dataset} reveals four critical physical properties of the PRONOSTIA dataset:
\begin{enumerate}[label=(\roman*),leftmargin=2.2em]
    \item \textbf{Extreme Lifetime Dispersion}: Under identical operating conditions, true lifespans vary by more than a factor of three (e.g., in Condition~1, \texttt{Bearing1\_2} failed at $8{,}710$~s, whereas \texttt{Bearing1\_1} survived to $28{,}030$~s).
    \item \textbf{Abrupt and Late Defect Onset}: The $3\sigma$ First Predicting Time (FPT) falls after 95\% of total life for 9 of the 17 bearings. For the majority of its operation, vibration is stationary Gaussian noise containing negligible prognostic information.
    \item \textbf{Inconsistent Failure Threshold}: Although the challenge protocol declared a $20$~g trip threshold, the terminal RMS at failure varies from $0.73$~g ($1.1\times$ baseline) to $14.25$~g ($31.7\times$ baseline). Fixed failure thresholds are therefore physically ungrounded.
    \item \textbf{Discontinuous Recording}: In \texttt{Bearing1\_1}, a $21{,}272$~s gap exists in the time stamps. In accordance with standard benchmark conventions, elapsed time is evaluated as snapshot indices multiplied by $10$~s.
\end{enumerate}

\section{Methodological Implementations}
\label{sec:annex_methods}

\subsection{Proposed KI-SSDM Pipeline for Bearings}
KI-SSDM maps the state-space degradation backbone onto rolling-element contact mechanics:
\begin{enumerate}[label=(\arabic*),leftmargin=2.2em]
    \item \textbf{Healthy Stage (Pre-Onset)}: Nominal fatigue life scales with ISO~281:2007~\cite{ISO_281} basic rating: $L_{10} \propto (C/P)^3 / n$. A hierarchical log-normal lifetime model conditions on survival to inspection time $t_{\mathrm{obs}}$:
    \begin{equation}
        \ln L_u = \mu_0 + \gamma g_c + \eta_u, \quad g_c = \ln\left[\left(\frac{P_1}{P_c}\right)^3 \left(\frac{n_1}{n_c}\right)\right], \quad \gamma \sim \mathcal{N}(1.0, 0.5^2), \; \eta_u \sim \mathcal{N}(0, \sigma_\eta^2),
    \end{equation}
    where $\mathbb{E}[T - t_{\mathrm{obs}} \mid T > t_{\mathrm{obs}}] = \int_{t_{\mathrm{obs}}}^\infty S(u)\,\mathrm{d}u / S(t_{\mathrm{obs}})$.
    \item \textbf{Defect Propagation Stage (Post-Onset)}: Following the $3\sigma$ FPT departure ($r_k > r_0 + 3\sigma_{\mathrm{base}}$), log-vibration defect growth $y_k = \ln(r_k / r_0)$ is tracked by a Kalman-GLS filter:
    \begin{equation}
        y_k = c_k + \beta_k \tau_k + B(\tau_k) + \varepsilon_k, \quad \mathrm{d}B(t) = -\theta_B B(t)\,\mathrm{d}t + \sigma_B\,\mathrm{d}W_B(t),
    \end{equation}
    where $B(t)$ is the Ornstein-Uhlenbeck nuisance process capturing reversible operational chatter and temperature fluctuations.
    \item \textbf{Smooth Bounded Learned Discrepancy}: Discrepancy multiplier $\alpha = \exp(\lambda \rho)$ is parameterised by a tanh-bounded ridge encoder ($|\rho| \le 1.0$) on causal feature descriptors (vibration ratio, kurtosis, spectral energy ratio, log-elapsed time). The pooling weight $\lambda \in \{0, 0.25, \dots, 1.0\}$ is selected via leave-one-out cross-validation on the training set, with $\lambda = 0$ as the pure physics fallback.
    \item \textbf{Stratified Split-Conformal UQ}: Nonconformity scores are calibrated separately for pre-onset and post-onset regimes on out-of-fold training predictions, ensuring rigorous finite-sample coverage guarantees.
\end{enumerate}

\subsection{Published Baseline Comparators}
\begin{enumerate}[label=(\alph*),leftmargin=2.2em]
    \item \textbf{Bayesian Exponential Model (Yan et al., 2021~\cite{Yan_2021})}: A conjugate linear Bayesian model tracking log-degradation signals with population-calibrated prior distributions. Before onset, it falls back to the median training life minus $t_{\mathrm{obs}}$.
    \item \textbf{Particle Filter Degradation Model (Deng et al., 2020~\cite{Deng_2020})}: An improved exponential degradation model $x_k = \phi + \theta \exp(\beta \tau_k)$ tracked by a Sequential Importance Resampling (SIR) particle filter with $3\sigma$ FPT detection and systematic resampling.
    \item \textbf{Deep LSTM Life-Percentage Network (Zhu et al., 2019~\cite{Zhu_2019}; Wang et al., 2023~\cite{Wang_2022})}: An authentic recurrent neural network with 30-snapshot rolling feature windows predicting normalized life percentage ($p = t / T_{\mathrm{fail}}$), with remaining life given by $\widehat{\RUL} = t(1-p)/p$.
    \item \textbf{Signal-Free Reference Baselines}: (1)~Life prior only (pure conditional survival without vibration features); (2)~Constant median RUL (predicting the static median RUL of training bearings).
\end{enumerate}

\section{Experimental Results and Statistical Analysis}
\label{sec:annex_results}

All methods were implemented and executed in MATLAB under strict causal protocols. No test-bearing labels, future data, or hyperparameter retuning was permitted.

\subsection{Protocol A: Official Challenge Split (11 Held-Out Bearings)}
Table~\ref{tab:pronostia_protocol_a} details performance under the official IEEE PHM 2012 challenge protocol (trained on 6 bearings, evaluated on 11 test bearings).

\begin{table*}[!t]
\centering
\caption{Prognostic performance on the 11 held-out test bearings of the IEEE PHM 2012 / PRONOSTIA dataset under Protocol~A (Official challenge split; Bearing1\_4 scored against official $339$~s).}
\label{tab:pronostia_protocol_a}
\small
\begin{tabular*}{\textwidth}{@{\extracolsep{\fill}}l c c c c c@{}}
\toprule
\textbf{Method} & \textbf{MAE (s)} & \textbf{PHM 2012 Score} & \textbf{90\% Coverage} & \textbf{Mean Width (s)} & \textbf{Late (\%)} \\
\midrule
\textbf{Proposed KI-SSDM} & \textbf{4,463} & \textbf{0.082} & \textbf{81.8\%} & \textbf{51,987} & 72.7\% \\
KI-SSDM (No Discrepancy) & 6,205 & 0.142 & 81.8\% & 213,346 & 81.8\% \\
KI-SSDM (No ISO Scaling) & 5,737 & 0.049 & 90.9\% & 220,826 & 90.9\% \\
KI-SSDM (No Survival Conditioning) & 5,500 & 0.018 & 90.9\% & 420,448 & 45.5\% \\
\midrule
Bayesian Exponential (Gebraeel 2005) & 6,224 & 0.128 & 81.8\% & 499,505 & 54.5\% \\
PF Exponential (Li 2015) & 7,300 & 0.089 & 90.9\% & 422,362 & 45.5\% \\
Deep LSTM Life-Percentage & 12,785 & 0.104 & 54.5\% & 119,885 & 81.8\% \\
\midrule
Reference: Life Prior Only (No Signal) & 4,648 & 0.129 & 90.9\% & 119,768 & 81.8\% \\
Reference: Constant Median RUL & 2,400 & 0.094 & 90.9\% & 22,587 & 72.7\% \\
\bottomrule
\end{tabular*}
\end{table*}

\subsection{Protocol B: Leave-One-Bearing-Out Cross-Validation (17 Bearings $\times$ 8 Origins)}
To eliminate small-sample partition bias, Table~\ref{tab:pronostia_protocol_b} reports leave-one-bearing-out (LOBO) evaluation across all 17 bearings evaluated at 8 origins each ($N=136$ total evaluations). Table~\ref{tab:pronostia_paired_tests} provides bearing-level paired statistical tests with family-wise Holm multiplicity adjustment across all 8 comparators.

\begin{table*}[!t]
\centering
\caption{Prognostic performance under Protocol~B (Leave-one-bearing-out cross-validation across all 17 bearings $\times$ 8 inspection origins, $N=136$ evaluations). Post-onset metrics evaluate the 39 origins occurring after $3\sigma$ FPT departure.}
\label{tab:pronostia_protocol_b}
\small
\begin{tabular*}{\textwidth}{@{\extracolsep{\fill}}l c c c c c@{}}
\toprule
\textbf{Method} & \textbf{MAE All (s)} & \textbf{MAE Post-Onset (s)} & \textbf{PHM Score} & \textbf{90\% Coverage} & \textbf{Mean Width (s)} \\
\midrule
\textbf{Proposed KI-SSDM} & \textbf{3,828} & \textbf{3,099} & \textbf{0.139} & \textbf{88.2\%} & \textbf{31,402} \\
KI-SSDM (No Discrepancy) & 5,611 & 4,971 & 0.098 & 89.7\% & 109,804 \\
KI-SSDM (No ISO Scaling) & 5,292 & 5,884 & 0.121 & 89.7\% & 88,118 \\
KI-SSDM (No Survival Conditioning) & 5,961 & 4,095 & 0.045 & 92.6\% & 425,147 \\
\midrule
Bayesian Exponential (Gebraeel 2005) & 6,696 & 6,661 & 0.040 & 90.4\% & 631,497 \\
PF Exponential (Li 2015) & 7,201 & 8,420 & 0.036 & 93.4\% & 436,490 \\
Deep LSTM Life-Percentage & 4,811 & 5,201 & 0.136 & 93.4\% & 63,334 \\
\midrule
Reference: Life Prior Only & 5,471 & 4,482 & 0.102 & 89.7\% & 109,070 \\
Reference: Constant Median RUL & 3,744 & 2,935 & 0.130 & 90.4\% & 22,634 \\
\bottomrule
\end{tabular*}
\end{table*}

\begin{table*}[!t]
\centering
\caption{Bearing-level paired statistical hypothesis tests comparing Proposed KI-SSDM against all comparators under Protocol~B ($n=17$ bearings, family-wise Holm multiplicity adjustment across 8 comparisons).}
\label{tab:pronostia_paired_tests}
\small
\begin{tabular*}{\textwidth}{@{\extracolsep{\fill}}l c c c c c@{}}
\toprule
\textbf{Comparator} & \textbf{Gain vs KI-SSDM [95\% CI] (s)} & \textbf{W / L} & \textbf{Holm $\bm{p_t}$} & \textbf{Holm $\bm{p_W}$} & \textbf{Statistical Verdict} \\
\midrule
Bayesian Exponential (Gebraeel 2005) & $+2{,}868$ [$+1{,}931, \; +3{,}805$] & 17 / 0 & $\mathbf{6.0 \times 10^{-5}}$ & $\mathbf{1.2 \times 10^{-4}}$ & KI-SSDM Significantly Better \\
PF Exponential (Li 2015)            & $+3{,}373$ [$+1{,}869, \; +4{,}876$] & 17 / 0 & $\mathbf{0.0013}$ & $\mathbf{1.2 \times 10^{-4}}$ & KI-SSDM Significantly Better \\
Deep LSTM Life-Percentage           & $+983$ [$-761, \; +2{,}727$]         & 11 / 6 & $0.50$ & $0.73$ & No Significant Difference (Tie) \\
Reference: Constant Median RUL       & $-84$ [$-490, \; +322$]              & 11 / 6 & $0.67$ & $0.85$ & No Significant Difference (Tie) \\
\midrule
Ablation: No Survival Conditioning  & $+2{,}132$ [$+1{,}344, \; +2{,}921$] & 16 / 1 & $\mathbf{0.0002}$ & $\mathbf{0.0003}$ & Survival Conditioning Essential \\
Ablation: No Discrepancy            & $+1{,}783$ [$+494, \; +3{,}071$]     & 10 / 7 & $\mathbf{0.049}$  & $0.25$ & Discrepancy Statistically Helpful \\
Ablation: No ISO Scaling            & $+1{,}463$ [$-340, \; +3{,}267$]     & 11 / 6 & $0.31$ & $0.73$ & Not Significant \\
\bottomrule
\end{tabular*}
\end{table*}

\begin{figure*}[!t]
\centering
\includegraphics[width=0.98\textwidth]{Fig13_PRONOSTIA_Benchmark.pdf}
\caption{Cross-domain benchmark validation on the IEEE PHM 2012 / PRONOSTIA bearing platform generated by \texttt{run\_pronostia\_benchmark.m}. (a--d)~RUL trajectories and 90\% split-conformal prediction intervals across operational life for representative bearings (\texttt{Bearing1\_3}, \texttt{Bearing1\_4}, \texttt{Bearing2\_3}, \texttt{Bearing3\_3}). (e)~Mean Absolute Error breakdown across all evaluation origins and post-onset origins. (f)~Empirical 90\% coverage alongside the official challenge asymmetric score. (g)~Bearing-level paired prognostic gain forest plot with 95\% confidence intervals. (h)~Systematic ablation across model components. Note: In figure panels, PI-DSSDM denotes the proposed KI-SSDM framework.}
\label{fig:pronostia_benchmark}
\end{figure*}

\section{Physical and Methodological Discussion}
\label{sec:annex_discussion}

Three primary insights emerge from the PRONOSTIA benchmark validation:
\begin{enumerate}[label=(\alph*),leftmargin=2.2em]
    \item \textbf{Significant Advantage over Classical Model-Based Prognostics}: KI-SSDM decisively outperforms the classical Bayesian exponential model (gain $+2{,}868$~s, $p < 10^{-4}$) and Particle Filter model (gain $+3{,}373$~s, $p = 0.0013$) across all 17 bearings (17/0). This advantage stems from two architectural elements: (1)~survival-conditioned prior modeling in the healthy stage (where Gebraeel and Li models are undefined and collapse to static subtractions); and (2)~causal feature tracking with bounded discrepancy estimation in Stage~II, preventing spurious vibration chatter from causing catastrophic gradient divergence.
    \item \textbf{Parity with Deep LSTM and Constant Median-RUL}: In paired statistical tests, KI-SSDM achieves parity with the Deep LSTM ($p=0.50$) and ties a signal-free constant median predictor ($p=0.67$). Under Protocol~A, the constant median predictor achieves an MAE of $2{,}400$~s, beating every algorithmic method. This occurs because the PRONOSTIA platform is characterized by extreme Weibull fatigue scatter under identical loading, accompanied by late, abrupt failure ($3\sigma$ departure after 95\% of life for 9 of the 17 bearings, with only a minority such as \texttt{Bearing1\_4} at 4.1\% and \texttt{Bearing2\_3} at 13.7\% exhibiting earlier onset). As a result, vibration signals provide virtually zero prognostic warning during the healthy operational phase.
    \item \textbf{Honesty of Split-Conformal Uncertainty Quantification}: Split-conformal calibration achieves $88.2\%$ empirical coverage under LOBO (close to the 90\% nominal target within finite-sample binomial bounds). However, the conformal prediction interval width ($\approx 31{,}400$~s, or $\approx 8.7$~hours) is wide. This proves that while conformal inference guarantees statistical validity without model-form assumptions, no mathematical wrapper can artificially reduce epistemic uncertainty when the underlying physical data exhibit three-fold lifetime variance without precursor telemetry signatures.
\end{enumerate}

\section{Software Implementation and Reproducibility}
\label{sec:annex_reproducibility}

All benchmark routines, raw feature extraction routines, and evaluation scripts are permanently archived in the accompanying package directory \texttt{pronostia/}:
\begin{itemize}[leftmargin=2.2em]
    \item \textbf{Feature Extraction}: \texttt{extract\_pronostia\_features.m} ingests raw FEMTO-ST accelerometer CSV files and exports cached features.
    \item \textbf{Benchmark Pipeline}: \texttt{run\_pronostia\_benchmark.m} executes both Protocol~A and Protocol~B, evaluating all 5 core models, 3 ablations, and 2 references across all 136 evaluation origins.
    \item \textbf{Execution}: Fully reproducible via headless batch command:
    \begin{center}
        \texttt{matlab -batch "cd pronostia; run\_pronostia\_benchmark();"}
    \end{center}
\end{itemize}

\begin{thebibliography}{99}
\bibitem{Lei_2018}
Y.~Lei, N.~Li, L.~Guo, N.~Li, T.~Yan, and J.~Lin, ``Machinery health prognostics: A systematic review from data acquisition to {RUL} prediction,'' \emph{Mechanical Systems and Signal Processing}, vol.~104, pp.~799--834, 2018.

\bibitem{ISO_281}
International Organization for Standardization, ``Rolling bearings --- Dynamic load ratings and rating life,'' \emph{ISO Standard 281:2007}, Geneva, Switzerland, 2007.

\bibitem{Yan_2021}
B.~Yan, X.~Ma, G.~Huang, and Y.~Zhao, ``Two-stage physics-based {Wiener} process models for online {RUL} prediction in field vibration data,'' \emph{Mechanical Systems and Signal Processing}, vol.~152, p.~107378, 2021.

\bibitem{Deng_2020}
Y.~Deng, S.~Du, S.~Jia, C.~Zhao, and Z.~Xie, ``Prognostic study of ball screws by ensemble data-driven particle filters,'' \emph{Journal of Manufacturing Systems}, vol.~56, pp.~359--372, 2020.

\bibitem{Zhu_2019}
J.~Zhu, N.~Chen, and W.~Peng, ``Estimation of bearing remaining useful life using multiscale convolutional neural network,'' \emph{IEEE Transactions on Industrial Electronics}, vol.~66, no.~4, pp.~3208--3216, 2019.

\bibitem{Wang_2022}
X.~Wang, L.~T.~Yang, E.~Cao, L.~Guo, L.~Ren, and M.~J.~Deen, ``A tensor-based {t-SVD-LSTM} remaining useful life prediction model for industrial intelligence,'' \emph{IEEE Transactions on Industrial Informatics}, vol.~19, no.~1, pp.~697--707, 2023.
\end{thebibliography}

\end{document}
